\section{Localization}

\subsection{Localization of rings}

Consider the following problem: for a ring $A$ and a subset $S\subseteq A$, we want a ring $S^{-1}A$ and a morphism $\iota\colon A\to S^{-1}A$
such that any morphism $f\colon A\to B$ such that $f(S)\subseteq B^\times$, $f$ factors uniquely through $\iota$.
That is, there is a unique $\hat f\colon S^{-1}A\to B$ such that $f=\hat f\circ\iota$:

\centerline{\drawdiagram{
    $A$&$S^{-1}A$\cr
    &$B$
}{
    \diagarrow{from={1,1}, to={1,2}, text=$\iota$, y distance=-.25cm}
    \diagarrow{from={1,1}, to={2,2}, text=$f$, y distance=-.25cm}
    \diagarrow{from={1,2}, to={2,2}, text=$\exists!$, x distance=.25cm, dashed}
}}

This is a universal property, and as such any solution is unique up to unique isomorphism.
In this section we will investigate the construction of such a solution.

\bdefn

    A subset $S\subseteq A$ is {\emphcolor multiplicatively closed} if $1\in S$ and for $a,b\in S$ we have $ab\in S$.

\edefn

It is sufficient to consider only the case that $S$ is multiplicatively closed.
Otherwise, let $\hat S$ be $S$'s multiplicative closure: the set of all finite products of elements from $S$.
Then $\hat S^{-1}R$ is a solution to the problem for $S$.
Indeed, if $f(S)\subseteq B^\times$, then $f(\hat S)\subseteq B^\times$ as well, as the product of units.
Then there is a unique factoring through $\hat S$'s $\iota$.

\bdefn

    For $S\subseteq A$ multiplicatively closed, define an equivalence on $A\times S$ by $(a,s)\sim(b,t)$ iff there is a
    $u\in S$ such that $u(at-bs)=0$, in which case we write $(a,s)\sim_u(b,t)$.

\edefn

\blemm

    This is indeed an equivalence relation.

\elemm

\bproof

    The subrelations $\sim_u$ are all clearly reflexive and symmetric, thus their union $\sim$ is as well.
    Now if $(a,s)\sim_\alpha(b,t)\sim_\beta(c,v)$ then $\alpha at=\alpha sb$ and $\beta bv=\beta ct$.
    So then $\alpha t\beta av=\alpha sb\beta v=\alpha\beta tcs$ so that $(a,s)\sim_{\alpha t\beta}(c,v)$.
    \qed

\eproof

\bexam

    A more naive equivalence would be $\sim_1$, i.e.\ $(a,s)\sim(b,t)$ iff $at-bs=0$.
    This corresponds to $\sim$ when $A$ is an integral domain, but is not generally an equivalence.

\eexam

\bdefn

    For $S\subseteq A$ multiplicatively closed, define $S^{-1}A$ to be the quotient of $A\times S$ by $\sim$.
    Write the equivalence class of $(a,s)$ by $a/s$.
    We give this a ring structure by
    $$ 0 = 0/1,\qquad 1 = 1/1,\qquad a/s + b/t = (at+bs)/st,\qquad a/s \cdot b/t = (ab)/(st) . $$
    Define $\iota\colon A\to S^{-1}A$ by $\iota(a)=a/1$.

\edefn

\bprop

    $S^{-1}A$ is indeed a ring, and $\iota$ is a ring morphism.

\eprop

\bproof

    We need to show that addition and multiplication are well-defined.
    Indeed, if $(a,s)\sim_u(a',s')$ and $(b,t)\sim_v(b',t')$ then $(at+bs,st)\sim_{uv}(a't'+b's',s't')$ which can be verified.
    And similarly $(ab,st)\sim_{uv}(a'b',s't')$.
    From well-definedness it is easy to see that addition and multiplication define a ring structure and $\iota$ is a ring morphism.
    \qed

\eproof

A few properties are immediate:
\benum
    \item For $s\in S$, $\iota(s)=1/s$ has an inverse $s/1$, so that $\iota(S)\subseteq(S^{-1}R)^\times$;
    \item We have that $a/s=0/1$ iff there is a $t\in S$ such that $at=0$.
    In particular, $\iota(a)=0$ iff there is a $t\in S$ such that $ta=0$.
    \item Every element of $S^{-1}A$ is of the form
    $$ a/s = a/1\cdot1/s = \iota(a)\cdot\iota(s)^{-1} . $$
\eenum

\bthrm

    This construction solves the above universal property.

\ethrm

\bproof

    Let $f\colon A\to B$ be such that $f(S)\subseteq B^\times$: we need to show that there is a unique map $g\colon S^{-1}A\to B$
    such that $f=g\circ\iota$.
    Indeed, we must define
    $$ g(a/s) = g(\iota(a)\iota(s)^{-1}) = (g\circ\iota(a))\cdot(g\circ\iota(s))^{-1} = f(a)f(s)^{-1} . $$
    This shows that $g$ is unique, if it exists.
    We need to show that this map is well-defined and a ring morphism.

    If $a/s=b/t$ then we must have that $f(a)f(s)^{-1}=f(b)f(t)^{-1}$ so that $g$ is well-defined.
    Indeed, if $u(at-bv)=0$ for $u\in S$, we have
    $$ f(u)\cdot(f(at)-f(bv)) = 0 $$
    and since $f(S)\subseteq B^\times$ we have that $f(u)$ is a unit so that $f(at)-f(bv)$, meaning this is well-defined.
    And this is a ring morphism, as can be checked easily.
    \qed

\eproof

\bcoro

    Let $f\colon A\to B$ be a ring morphism such that
    \benum
        \item $f(S)\subseteq B^\times$,
        \item if $f(a)=0$ then $ta=0$ for some $t\in S$,
        \item every element of $B$ is of the form $f(a)f(s)^{-1}$ for $a\in A$ and $s\in S$.
    \eenum
    Then there is a unique isomorphism $g\colon S^{-1}A\to B$ such that $g\circ\iota=f$.

\ecoro

\bproof

    $f$ solves the universal problem above for $S$ as well.
    If $h\colon A\to C$ is a map where $h(S)\subseteq C^\times$, then we need to define $\hat h\colon B\to C$ by
    $$ \hat h(f(a)f(s)^{-1}) = h(a)h(s)^{-1} , $$
    which is well-defined and a ring morphism for similar reasons as above.
    Since universal properties are unique up to unique isomorphism, the result follows.

    Alternatively we can prove this directly: let $g\colon S^{-1}A\to B$ be the unique map where $g\circ\iota=f$ by the universal property.
    By definition $g(a/s)=f(a)f(s^{-1})$.
    By the third point $g$ is surjective, and it is injective since if $g(a/s)=0$ we have $f(a)=0$ so that $ta=0$, which means that $a/s=0$.
    \qed

\eproof

\bexam

    If $A$ is an integral domain, then $S=A-0$ is multiplicatively closed, and $S^{-1}A$ is the field of fractions of $A$.

\eexam

\bprop

    For $A\neq0$, then $S^{-1}A=0$ iff $0\in S$.

\eprop

\bproof

    We have that $S^{-1}A=0$ iff $1/1=0/1$ iff there is a $s\in S$ such that $s\cdot1=0$.
    \qed

\eproof

\bdefn

    Let $\p\subseteq A$ be prime.
    Then $S=A-\p$ is multiplicatively closed, and we denote its localization by $A_\p=S^{-1}A$.

    For $a\in A$, $S=\set{1,a,a^2,\dots}$ is multiplicatively closed, and we denote its localization by $A_a=S^{-1}A$.

\edefn

For $a\in A$ not nilpotent (so $A_a\neq0$) define a map
$$ \phi\colon A[x] \to A_a,\qquad x\mapsto 1/a . $$
This map is clearly surjective, and $p(x)=\sum_k^nb_kx^k\in A[x]$ is mapped to zero iff
$$ \sum_{k=0}^nb_k\frac1{a^k} = \frac{\sum_kb_ka^{n-k}}{a^n} = 0 . $$
This is iff there is a $t$ such that
$$ a^t\sum_kb_ka^{n-k} = 0 . $$
Defining $c_k$ recursively by $c_0=-b_0$ and $c_k=ac_{k-1}-b_k$ gives a polynomial such that $p(x)=(ax-1)q(x)$, so $p\in(ax-1)$
(the equality above ensures the recursion terminates).
Thus $\ker\phi=(ax-1)$ and thus
$$ A_f \cong \frac{A[x]}{(ax-1)} . $$

\bnote

    Note that we have two distinct localizations here, which have similar names in some cases.
    If $A$ is a PID and $p\in A$ is irreducible, then $A_{(p)}$ is localization by $A-(p)$, which $A_p$ is localization by $\set{p^n}_{n<\omega}$.
    Elements of the former are of the form $a/s$ where $p$ does not divide $s$, while elements of the former are of the form $a/s$ where $s=p^n$.
    Quite different!

\enote

\bdefn

    For a ring morphism $f\colon A\to B$ and an ideal $I\subseteq A$, define the ideal $IB=(f(I))$ of $B$.
    Explicitly, $IB$ is the set of all finite sums $\sum_if(a_i)b_i$ for $a_i\in I$.

\edefn

In general we have $f^{-1}(f(I))\supseteq I$ and $f(f^{-1}(J))\subseteq J$.
These imply $f^{-1}(IB)\supseteq I$ and $f^{-1}(J)B\subseteq J$.

\bdefn

    For two ideals $I,J$, define $(I:J)=\set{x\in A}[xJ\subseteq I]$.
    This is an ideal containing $I$.
    For $a\in A$, define $(I:a)=(I:(a))=\set{x\in A}[xa\subseteq I]$.

    For $S\subseteq R$ multiplicatively closed, define the ideal $S^{-1}I=IS^{-1}A\subseteq S^{-1}A$.
    This is the ideal generated by $\iota(I)$.

\edefn

So we have two maps: $S^{-1}\bul$ mapping ideals of $A$ to ideals of $S^{-1}A$, and $\iota^{-1}$ mapping ideals of $S^{-1}A$ to $A$:
$$ \set{\hbox{Ideals of $A$}} \xtof{\quad S^{-1}\quad}[\quad\iota^{-1}\quad] \set{\hbox{Ideals of $S^{-1}A$}} . $$
We wish to relate the two.

\bprop

    \benum
        \item For an ideal $I\subseteq A$, $S^{-1}I=\set{a/s}[a\in I,s\in S]$.
        \item For an ideal $J\subseteq S^{-1}A$, we have $S^{-1}\iota^{-1}J=J$.
        \item For an ideal $I\subseteq A$, we have $\iota^{-1}S^{-1}I=\bigcup_{s\in S}(I:s)$.
        \item We have $S^{-1}I=S^{-1}A$ iff $I\cap S\neq\emptyset$.
        \item For $\p\subseteq A$ prime, if $I\cap S\neq\emptyset$ then $S^{-1}\p\subseteq S^{-1}A$ is prime (otherwise it is all of $S^{-1}A$).
        \item The correspondence $\q\mapsto\iota^{-1}\q$ is a bijection between prime ideals of $S^{-1}A$ and prime ideals of $A$ not intersecting $S$.
        Its inverse is $S^{-1}$.
        That is, we have a one-to-one correspondence:
        $$ \set{\hbox{Prime ideals of $A$ not intersecting $S$}} \xtof{\quad S^{-1}\quad}[\quad\iota^{-1}\quad] \set{\hbox{Prime ideals of $S^{-1}A$}} . $$
    \eenum

\eprop

\bproof

    \benum
        \item $\set{a/s}[a\in I,s\in S]$ is an ideal and contains $\iota(I)$.
        It is clearly closed under multiplication by $S^{-1}A$ and for $a/s+b/t=(at+bs)/(st)$, notice that $at+bs\in I$.
        By definition $S^{-1}I$ is the smallest ideal containing $\iota(I)$, and it consists of finite sums $\sum_i\iota(a_i)x_i/s_i$.
        This clearly contains $a/s$, and thus they must be equal.

        \item As explained before, in general we have $S^{-1}\iota^{-1}J=\iota^{-1}(J)S^{-1}A\subseteq J$, we must prove the opposite inclusion.
        As above, we have $S^{-1}\iota^{-1}J=\set{a/s}[\iota(a)\in J,s\in S]$.
        For $x/s\in J$, we have that $\iota(x)=x/1\in J$ and so $x/s\in S^{-1}\iota^{-1}J$ as desired.

        \item If $a\in\iota^{-1}S^{-1}I$ then $\iota(a)\in S^{-1}I$ so by the first point $a/1=b/s$ for $b\in I$.
        Then $asu=bu$ for some $u\in S$ and since $bu\in J$ this means that $a\in(I:su)$.
        Conversely if $a\in(I:s)$ for $s\in S$ then $as\in I$ and so $\iota(a)=a/1=(as)/s\in S^{-1}I$.

        \item If $S^{-1}I=S^{-1}A$ then $1/s=a/t$ for $a\in I$, so $as=t$ meaning $as\in I\cap S$.
        Conversely if $a\in I\cap S$ then $1/1=a/a\in S^{-1}I$.

        \item If $S\cap\p=\emptyset$ then we need to show that $S^{-1}\p$ is prime.
        If $a_1/s_1,a_2/s_2\notin S^{-1}\p$ we want to show $(a_1a_2)/(s_1s_2)\notin S^{-1}\p$.
        By assumption $a_1,a_2\notin\p$.

        Now we claim that if $a\notin\p$, $a/s\notin S^{-1}\p$, which would complete the proof.
        Assume the contrary: $a/s=b/t$ for $b\in\p$ then there is a $u\in S$ such that $atu=bsu$.
        Then $atu\in\p$, but $a\notin\p$ and $t,u\in S$ and thus also are not in $\p$, a contradiction to $\p$'s primality.

        \item Preimages of prime ideals are prime, so $\iota^{-1}$ maps prime ideals to prime ideals.
        By the second point $S^{-1}\iota^{-1}\q=\q$ for all prime ideals $\q\subseteq S^{-1}A$.
        It remains to show that $\iota^{-1}S^{-1}\p=\p$ for all prime $\p\subseteq A$ which don't intersect $S$.
        We know that $\p\subseteq\iota^{-1}S^{-1}\p$ as is true in general, so we want to show that $\iota^{-1}S^{-1}\p=\bigcup_{s\in S}(\p:s)$ is contained in $\p$.
        Indeed, if $a\in(\p:s)$ then $as\in\p$ and since $s\notin\p$ we have that $a\in\p$ as desired.
        \qed
    \eenum

\eproof

\bcoro

    Let $\p\subseteq A$ be prime, and consider its localization $A_\p$.
    Then $A_\p$ is a local ring with maximal ideal $\p A_\p$.
    Moreover, the correspondence $\q\to\q A_\p$ is a bijection between prime ideals of $A_\p$ and prime ideals $\q\subseteq\p$.

\ecoro

\bproof

    Let $S=A-\p$, so that $A_\p=S^{-1}A$, and $\p A_\p=S^{-1}\p$.
    The last point is clear: $\q\cap S=\emptyset$ iff $\q\subseteq\p$.
    $A_\p$ being local is also clear: $S^{-1}\bul$ is order-preserving, so $S^{-1}\q\subseteq S^{-1}\p$ for all prime ideals $\q$.
    Thus $S^{-1}\p$ is maximal.
    \qed

\eproof

\bcoro

    If $a\in A$ is not nilpotent, then there is a prime ideal $\p\subseteq A$ such that $a\notin\p$.

\ecoro

\bproof

    As before, consider the multiplicative set $S=\set{a^n}_{n<\omega}$, and $A_a=S^{-1}A$.
    Since $a$ is not nilpotent, $0\notin S$, so that $A_a\neq0$.
    So take a prime ideal $\p\subseteq A_a$, and then $\iota^{-1}\p$ is a prime ideal of $A$ not intersecting $S$, thus not containing $a$.
    \qed

\eproof

\bnote

    This completes the proof before that
    $$ \nil(A) = \bigcap_\p\p $$
    but without appealing to Zorn, and thus not relying on the Axiom of Choice.

\enote

\subsection{Localization of modules}

\bdefn

    Let $S\subseteq A$ be multiplicatively closed and $M$ an $A$-module.
    Define an equivalence on $M\times S$ by $(m,s)\sim(m',s')$ iff there is a $t\in S$ such that $t(s'm-sm')=0$.
    This is an equivalence relation and we denote the equivalence class of $(m,s)$ by $m/s$.
    The quotient $S^{-1}M$ forms an $A$-module by
    $$ m/s + m'/s' = (s'm+sm')/(ss'),\qquad a\cdot(m/s) = (am)/s . $$

\edefn

Note that localization at $S$ forms an additive functor $S^{-1}\colon\Mod_A\to\Mod_{S^{-1}A}$.
Indeed, $S^{-1}M$ is an $S^{-1}A$ module by
$$ a/s\cdot m/s' = (am)/(ss') . $$
And if $f\colon M\to N$ is an $A$-module morphism, then define
$$ S^{-1}f\colon S^{-1}M\to S^{-1}N,\qquad S^{-1}f(m/s) = f(m)/s . $$
It is not hard to see that this is functorial and additive.

There is the canonical morphism $\iota\colon A\to S^{-1}A$ and as such it lifts to a ``forgetful'' functor $\Mod_{S^{-1}A}\to\Mod_A$,
where for $M\in\Mod_{S^{-1}A}$ we view it as an $A$-module by $a\cdot M=\iota(a)\cdot M$.

\bprop

    The functor $S^{-1}\colon\Mod_A\to\Mod_{S^{-1}A}$ is left adjoint to the forgetful functor $\Mod_{S^{-1}A}\to\Mod_A$.

\eprop

\bproof

    We need to show that there is a canonical bijection
    $$ \hom_{S^{-1}A}(S^{-1}M,N) \cong \hom_A(M,N) , $$
    for all $A$-modules $M$ and $S^{-1}A$-modules $N$.
    For $f\colon M\to N$ map it to $\hat f\colon S^{-1}M\to N$ by $\hat f(m/s)=1/s\cdot f(m)$.
    And for $g\colon S^{-1}M\to N$ map it to $\tilde g\colon M\to N$ by $\tilde g(m)=g(m/1)$.
    \qed

\eproof

\bdefn

    An additive functor $F\colon\Mod_A\to\Mod_B$ is
    \benum
        \item {\emphcolor left exact} if it preserves left-exact sequences: if
        $$ 0 \to M \xtos f N \xtos g P $$
        is exact, then so is
        $$ 0 \to FM \xtos{Ff} FN \xtos{Fg} FP . $$
        \item {\emphcolor right exact} if it preserves right-exact sequences: if
        $$ M \xtos f N \xtos g P \to 0 $$
        is exact, then so is
        $$ FM \xtos{Ff} FN \xtos{Fg} FP\to 0. $$
        \item {\emphcolor exact} if it is both left and right exact.
    \eenum

\edefn

\bprop

    An additive functor is
    \benum
        \item left exact iff it preserves kernels iff it preserves finite limits iff it preserves the left side of a short exact sequence.
        \item right exact iff it preserves cokernels iff it preserves finite colimits iff it preserves the right side of a short exact sequence.
        \item exact iff it preserves finite limits and colimits iff it preserves short exact sequences iff it preserves all exact sequences.
    \eenum

\eprop

\bproof

    Left-exact sequences are kernel diagrams, and right-exact sequences are cokernel diagrams.
    So by definition, left exact functors are those which preserve kernels, and right exact functors are those which preserve cokernels.
    Since additive functors preserve finite biproducts, left/right exact functors preserve limits/colimits.

    Clearly left-exact functors preserve the left side of a short exact sequence (similarly for right-exact functors).
    To show the converse, if left-exact sequences can be quotiented to give an appropriate short exact sequence, similar for right-exact
    sequences.
    \qed

\eproof

\bcoro

    Left adjoints are right exact, and right adjoints are left exact.

\ecoro

\bproof

    Left adjoints preserve all colimits, and right adjoints all limits.
    \qed

\eproof

\bcoro

    $S^{-1}$ is an exact functor.

\ecoro

\bproof

    By the above corollary, it is right-exact.
    To show that it is exact, it is sufficient to show that it preserves injections.
    If $f\colon M\to N$ is injective then $m/s\in\ker S^{-1}f$ iff $f(m)/s=0/1$, so that $tf(m)=0$ for some $t\in S$.
    But then $tf(m)=f(tm)$ so $tm\in\ker\phi$, so $tm=0$, which means $m/s=0/1$ as desired.
    \qed

\eproof

In general if $U\colon\Mod_A\to\Mod_B$ is a restriction functor corresponding to a morphism $\iota\colon B\to A$, then it has a left adjoint
via the tensor product.
In particular, viewing $A$ as a $(A,B)$-module, then $A\otimes_BM$ has an $A$-module structure (given by $a(a'\otimes m)=(aa')\otimes m$).
$U$'s left adjoint is then $A\otimes_BM$.

Since left adjoints are unique up to natural isomorphism,
$$ S^{-1}M \cong S^{-1}A\otimes_AM . $$
Thus $S^{-1}A\otimes_A\bul$ is an exact functor, which means that $S^{-1}A$ is a {\emphcolor flat} $A$-module.
This can also be shown directly: $S^{-1}A$ is the filtered colimit of free $A$-modules, and is thus flat.

\bcoro

    $S^{-1}$ preserves arbitrary coproducts:
    $$ S^{-1}\bigoplus_\alpha M_\alpha \cong \bigoplus_\alpha S^{-1}M_\alpha $$
    naturally.

\ecoro

\bproof

    Left adjoints preserve arbitrary colimits.
    \qed

\eproof

\bnote

    While $S^{-1}$ is exact, it is not {\emphcolor continuous}: it does not preserve arbitrary limits (equivalently it does not preserve
    arbitrary products).
    That is, we do {\it not\/} have that
    $$ S^{-1}\prod_\alpha M_\alpha \cong \prod_\alpha S^{-1}M_\alpha . $$
    This is of course true when the index set is finite, since then it is just a biproduct.

\enote

\bcoro

    \benum
        \item If $N\subseteq M$ is a submodule, $S^{-1}N$ can be naturally embedded in $S^{-1}M$.
        \item There is a natural isomorphism $S^{-1}(M/N)\cong (S^{-1}M)/(S^{-1}N)$.
        \item For $f\colon M\to N$ of $A$-modules, we have natural isomorphisms
        $$ \ker S^{-1}f \cong S^{-1}\ker f,\qquad \im S^{-1}f \cong S^{-1}\im f,\qquad S^{-1}\coker f\cong\coker S^{-1}f . $$
        \item For every two submodules $M_1,M_2\subseteq M$ under the natural embeddings,
        $$ S^{-1}(M_1\cap M_2) = S^{-1}M_1\cap S^{-1}M_2 \subseteq S^{-1}M . $$
    \eenum

\ecoro

\bproof

    The first point is due to $S^{-1}$ being exact and thus preserving injections.
    The second point is because $S^{-1}$ preserves the short exact sequence $0\to N\to M\to M/N\to0$.
    The third point is clear as well, with only $\im$ not being immediately obvious.
    $\im f$ can be characterized as $\im f=\ker\coker f$, so it follows.

    For the last point, note that $M_1\cap M_2$ is the kernel of $M_1\oplus M_2\to M$ by $(x,y)\mapsto x-y$
    (via the diagonal embedding), and so this follows from $S^{-1}$'s exactness.
    \qed

\eproof

\bdefn

    For a prime ideal $\p\leq A$ and an $A$-module $M$, denote by $M_\p$ the localization of $M$ by $S=A-\p$.

\edefn

\blemm

    For an $A$-module $M$, consider the natural map
    $$ \phi\colon M\to\prod_\m M_\m,\qquad m\mapsto (m/1)_\m , $$
    where the product is over all maximal ideals of $A$.
    Then $\phi$ is injective.

\elemm

\bproof

    Assume there is an $0\neq x\in M$ such that $\phi(x)=0$.
    Then define
    $$ I = \ann_A(x) = \set{a\in A}[ax=0] . $$
    $I$ is a proper ideal of $A$, and thus is contained in a maximal ideal $\m$.
    Then for every $s\in A-\m$ we have that $s\notin I$ and so $sx\neq0$, and thus $m/1$ is not zero in $M_\m$, a contradiction.
    \qed

\eproof

\bcoro

    For an $A$-module $M$, the following are equivalent:
    \benum
        \item $M=0$,
        \item $S^{-1}M=0$ for every multiplicative set $S\subseteq A$,
        \item $M_\p=0$ for every prime ideal $\p\leq A$,
        \item $M_\m=0$ for every maximal ideal $\m\leq A$.
    \eenum

\ecoro

\bproof

    Clearly each point implies the next.
    By above $M$ embeds into $\prod_\m M_\m$, so if $M_\m=0$ for all $\m$ then $M=0$; the last point implies the first.
    \qed

\eproof

\bcoro

    For a morphism of $A$-modules $f\colon M\to N$, the following are equivalent:
    \benum
        \item $f$ is injective/surjective/an isomorphism,
        \item $S^{-1}f$ is injective/surjective/an isomorphism for every multiplicative $S\subseteq A$,
        \item $f_\p$ is injective/surjective/an isomorphism for every prime $\p\leq A$,
        \item $f_\m$ is injective/surjective/an isomorphism for every maximal $\m\leq A$.
    \eenum

\ecoro

\bproof

    By $S^{-1}$'s exactness, the first point implies the second, and clearly all the other points imply the next.
    To show the last implies the first, let $K=\ker f$ so that $0\to K\to M\xto fN$ is exact.
    Then $0\to K_\m\to M_\m\xto{f_\m}N_\m$ is exact, and so $K_\m=\ker f_\m=0$ for all $\m$.
    By above this means $K=0$ so that $f$ is injective.
    Similar for surjectivity.
    \qed

\eproof

